% Master abstract. Journal-specific main files override or wrap this.
\noindent
Casualty figures in armed conflict are reported by sources with opposing incentives: belligerents inflate enemy combatants killed and deflate civilians; advocacy and intergovernmental bodies push the other way. We treat the civilian-to-combatant death ratio as a partially identified parameter and ask what its value can be when conflicting sources are aggregated. The paper delivers two objects.

\textbf{(1) Bounds from under-fudgible relations.} Under a one-parameter model in which civilian adult men face $\mu$ times the death rate of women and children, we give a sharp identified set for the combatant share $q=D_M/D$ as a function of population structure, the demographic composition of the dead, and pre-war combatant stock; intersected with the manpower bound $D_M\le M$, this is a closed-form interval. Because the inputs are measured by different, non-aligned institutions, the relation is under-fudgible: moving any one input to rescue a disputed claim forces a compensating, visible move in an independently measured quantity. We define the \emph{contradiction radius} $\rho$ of a claim---the smallest standardised revision of a single independent input that makes it feasible---turning ``we dispute your number'' into ``name which measurement is wrong, and by how much.''

\textbf{(2) Sensitivity bands under political bias.} Each source is allowed to be Huber $\varepsilon_i$-contaminated; the analyst's political-bias prior on a source enters as $\varepsilon_i$. We give a closed-form first-order sensitivity band that absorbs the assumed contamination, a patternwise worst-case envelope, and a quantile-displacement bound under an explicit posterior-weight condition.

In the Gaza 2023--2026 application---anchored on the Ministry of Health demographic record, whose composition is consistent with an independent household survey---the identified set for $q$ is $[0,\nQOneRound{}\%]$ under only the weak assumption that civilian men are at least as exposed as women and children ($\mu\ge 1$), and $[0,\nQTwoRound{}\%]$ under exposure ratios informed by historical conflicts with known combatant status ($\mu\in[2,3.5]$). The public Israel Defense Forces figure of $17$--$25$k combatants killed converts to $q\approx \nIdfQLoRound{}$--$\nIdfQHiRound{}\%$ over the recorded toll. Its upper end is arithmetically infeasible: no exposure ratio $\mu\ge 1$ admits it, at a contradiction radius of $\nRhoOmegaAgRound{}$--$\nRhoOhchrCalRound{}$ times the conservative uncertainty scale attached to the demographic record. Its lower end ($q\approx \nIdfQLoRound{}\%$) sits at the very edge of the $\mu\ge 1$ set---attainable only if civilian men died at essentially the same rate as women and children ($\mu\le \nMuIdfLo{}$)---and leaves the set beyond that point, at a radius of $\nRhoIdfLoCalRound{}$ such units against the historically informed range. The rejection under calibrated exposure survives an upward correction of the death toll and, separately, a $5\%$ contamination allowance on each core input in the paper's first-order sensitivity band. Independent evidence---demographic decompositions, the military's own leaked internal database of named militant dead, and a fully documented spatial Bayesian model---concentrates well below the claim, at civilian-to-combatant ratios of roughly $\nCcrAtQOne{}{:}1$ to $\nSpatialCcr{}{:}1$ on direct deaths. The framework is symmetric: it tests the polar claim of zero combatant deaths the same way; that claim survives the demographic test and is rejected on documentary grounds.
